\begin{table}[ht]
\caption{Balance Table of office attributes by drafted status}
\centering
\begin{tabular}{lrrrr}
  \hline
 Year: 1931 & Not Drafted & Drafted & Difference & T‑stat \\ 
  \hline
  Female \%        & 0.02 & 0.02 & 0.00 & -0.33 \\ 
  Civil Servant \% & 0.09 & 0.06 & -0.03 & 1.01 \\ 
  \# of members         & 35.25 & 108.67 & 73.42 & -3.76 \\ 
  \hline 
\end{tabular}
\end{table}

\begin{table}[ht]
\centering
\begin{tabular}{lrrrr}
  \hline
 Year: 1936 & Not Drafted & Drafted & Difference & T‑stat \\ 
  \hline
  Female \%        & 0.04 & 0.02 & -0.02 & 3.07 \\ 
  New Hire \%      & 0.49 & 0.53 & 0.04 & -1.09 \\ 
  Civil Servant \% & 0.06 & 0.09 & 0.03 & -1.41 \\ 
  \# of members    & 18.02 & 56.48 & 38.46 & -4.96 \\ 
  \hline 
\end{tabular}
\end{table}
\note{Here I'm showing a balance table of drafted and non-drafted workers. My main points is that there isn't significant difference in experience with working with women. Both years show small numbers, which underlines how gender balance was skewed in the Tokyo civil service. Notably, I find that larger offices have higher chances of getting drafted. This is concerning factor as office size correlates with drafting and female shares.}